\documentclass[11pt,a4paper]{article}

% ---------- Codifica e lingua ----------
\usepackage[T1]{fontenc}
\usepackage[utf8]{inputenc}
\usepackage[main=italian,provide=*]{babel}

% ---------- Matematica ----------
\usepackage{amsmath,amssymb}
\usepackage{mathtools}

% ---------- Impaginazione ----------
\usepackage[a4paper,margin=2.7cm]{geometry}
\usepackage{caption}
\usepackage{booktabs}
\usepackage{enumitem}
\usepackage{placeins}
\usepackage[colorlinks=true,linkcolor=black,citecolor=black,urlcolor=blue!55!black]{hyperref}

\newcommand{\ket}[1]{\lvert #1 \rangle}
\newcommand{\bra}[1]{\langle #1 \rvert}

\title{Rumore quantistico sul trimero ad anello\\
\large Panoramica della pipeline (vista a 2000 piedi)}
\author{Samuele \\ Relatore: Prof.\ Alessandro Chiesa}
\date{}

\begin{document}
\maketitle

\begin{abstract}
\noindent Questo documento è una mappa, non una derivazione: serve a
vedere in un colpo d'occhio come è strutturata la pipeline di rumore sul
trimero ad anello prima di entrare nel dettaglio di ciascuno stadio
(ciascuno ha un proprio documento dedicato, elencato in
Sez.~\ref{sec:stadi}). È il mirror esatto della pipeline già completata
sul dimero (\texttt{panoramica\_\allowbreak pipeline\_rumore\_\allowbreak dimero.tex}): stessi
cinque stadi, stesso modello di rumore, stessa terminologia. Cambia il
sistema --- da 2 a 3 spin, ansatz $W$-2q.6 al posto di PMA-2q$\cdot$3 ---
non la struttura della pipeline. Scope e canali di rumore restano quelli
confermati dal relatore (\texttt{domande\_relatore.md}, sez.~12) e
topologia confermata (sez. ``topologia del trimero per il rumore'',
6 settembre 2026): l'\textbf{anello}, non la catena.
\end{abstract}

\tableofcontents

\section{Da dove partiamo}

La parte a sistema chiuso è completa sul trimero ad anello: teoria
esatta (decomposizione di Kambe, termine DM Opzione B), VQE senza e con
DM (ansatz $W$-2q.6), evoluzione temporale $e^{-iHt}$ via
Suzuki--Trotter, e correlazioni dinamiche $C_{ij}^{\alpha\beta}(t)$
tramite lo stesso circuito di tipo Hadamard test usato sul dimero. Tutto
lavora su uno stato \emph{puro} $\ket{\psi}$ a 3 qubit che evolve
unitariamente: non c'è interazione con l'ambiente, non c'è rumore.

Questa pipeline ripercorre esattamente la stessa catena di stadi già
fatta sul dimero --- preparazione dello stato fondamentale, evoluzione
temporale, misura dei correlatori --- ma la accompagna con le stesse due
sorgenti di rumore realistiche già calibrate per il dimero. Non si
introduce fisica nuova nell'Hamiltoniana né un nuovo modello di rumore:
si applica lo stesso modello, già validato, a un circuito a 3 qubit
invece che 2.

\section{Cosa cambia: da \texorpdfstring{$\ket{\psi}$}{|psi>} a \texorpdfstring{$\rho$}{rho}}

Come per il dimero, il rumore accoppia il sistema a un ambiente che non
osserviamo: uno stato puro non basta più, serve l'operatore densità
$\rho$ e un'evoluzione non unitaria pura. In Qiskit questo corrisponde,
esattamente come per il dimero, a usare
\texttt{AerSimulator(method=\allowbreak"density\_matrix")} invece del
simulatore statevector --- qui su un registro a 3 qubit.

\section{Gli stadi della pipeline}
\label{sec:stadi}

\begin{table}[htbp]
\centering
\begin{tabular}{@{}cll@{}}
\toprule
\# & Stadio & Cosa cambia rispetto al sistema chiuso \\
\midrule
1 & Modello di rumore & wrapper sul dimero, stessa logica \\
2 & VQE con termine DM, rumoroso & ansatz $W$-2q.6, simulatore a $\rho$ \\
3 & Quantum simulation (Trotter), rumorosa & stesso circuito, compilato per passo \\
4 & Correlazioni dinamiche, rumorose & stesso Hadamard test, readout incluso \\
5 & Scan sui parametri di rumore & nuovo: nessun analogo nel sistema chiuso \\
\bottomrule
\end{tabular}
\caption{I cinque stadi della pipeline, mirror esatto di quelli del
dimero. Il primo stadio non richiede alcun adattamento della
\emph{fisica} del modello: il \texttt{NoiseModel} non dipende dal
numero di qubit per costruzione (verificato, non assunto ---
Sez.~\ref{subsec:modello}). Un nuovo file esiste comunque
(\texttt{noise\_model\_\allowbreak trimero\_anello.py}), ma è un
wrapper sottile, non una reimplementazione.}
\end{table}
\FloatBarrier

\subsection{Modello di rumore}
\label{subsec:modello}

Stessi due canali confermati dal relatore per il dimero (errore di gate
depolarizzante 1q/2q, errore di lettura simmetrico), stessa calibrazione
reale (\texttt{ibm\_torino}, arXiv:2504.15187). La logica del
\texttt{NoiseModel} non viene reimplementata: costruisce l'errore con
\texttt{add\_all\_qubit\_\allowbreak quantum\_error}/\texttt{add\_all\_qubit\_\allowbreak readout\_error},
operazioni intrinsecamente indipendenti dal numero di qubit del circuito
su cui vengono applicate. \texttt{noise\_model\_\allowbreak trimero\_anello.py} è un
\textbf{wrapper sottile} che importa e ri-espone questa logica da
\texttt{noise\_model\_dimero.py} (dipendenza dichiarata esplicitamente
nel docstring, non nascosta), cosi' che il trimero abbia un punto
d'ingresso con lo stesso nome usato in ogni altra parte del progetto,
senza duplicare codice da tenere sincronizzato. Verificato empiricamente
sul circuito VQE+DM a 3 qubit, non solo argomentato dalla forma del
codice. Documento:
\texttt{risultati\_\allowbreak modello\_rumore\_\allowbreak trimero\_anello.tex}.

\subsection{VQE con termine DM, sotto rumore}

Si ripete la stima del ground state con l'ansatz $W$-2q.6 già validato a
sistema chiuso (fedeltà $\approx1$ nel caso ideale, punto di lavoro
$J=1,\,J'=0.4,\,b=b_c=2.4,\,D=0.15$), eseguito su simulatore a operatore
densità con il \texttt{NoiseModel} agganciato. Documento:
\texttt{risultati\_vqe\_dm\_\allowbreak rumoroso\_trimero\_anello.tex}.

\subsection{Quantum simulation (Trotter), sotto rumore}
\label{subsec:trotter}

Stesso circuito di Suzuki--Trotter del sistema chiuso, stessa insidia
già affrontata sul dimero: transpilare l'\emph{intero} circuito a un
livello di ottimizzazione alto collasserebbe gli $N$ passi identici in
un costo costante, cancellando la dipendenza da $N$ che serve allo scan.
A differenza del dimero, qui la strategia di transpilazione a blocchi
va scritta da zero: né
\texttt{circuito\_\allowbreak correlazioni\_\allowbreak trimero\_anello.py}
né \texttt{trotter\_trimero\_anello.py} la contengono già (verificato,
non assunto). Soluzione: transpilare una volta il singolo passo/blocco
al suo costo minimo, poi ripetere il blocco già compilato $N$ volte.

\textbf{Due scenari trattati}, non uno solo: Parte 1 dell'anello usa
punti di lavoro diversi per il VQE ($b=b_c=2.4,\,D=0.15$) e per la
dimostrazione Trotter ($R_0$: $b=0.05,\,D=1.93$, provvisorio, da
$\ket{000}$) --- a differenza del dimero, dove un unico punto serve
entrambi. Lo Stadio 3 copre entrambi: scenario A (preparazione VQE + $N$
passi) mirror esatto del dimero, scenario B (nessuna preparazione, da
$\ket{000}$) mirror della dimostrazione di Parte 1, reso rumoroso.
Documento:
\texttt{risultati\_\allowbreak trotter\_\allowbreak rumoroso\_trimero\_anello.tex}.

\subsection{Correlazioni dinamiche, sotto rumore}

Stesso circuito Hadamard test del sistema chiuso (ancilla + 3 qubit di
registro), ora con rumore su preparazione, blocco $U(t)$ ripetuto e
misura sull'ancilla (readout incluso). Documento:
\texttt{risultati\_correlatori\_\allowbreak rumorosi\_trimero\_anello.tex}.

\subsection{Scan sui parametri di rumore}

Come per il dimero: si ripete il conto dello stadio precedente al
variare di $(\varepsilon_{1q},\varepsilon_{2q},p_\text{readout})$, per
vedere come si sposta il numero di passi ottimale $N^*$. Documento:
\texttt{risultati\_scan\_\allowbreak parametri\_rumore\_trimero\_anello.tex}.

\section{Cosa resta fuori}

Stesso scope del dimero, per coerenza dichiarata esplicitamente:
$T_1/T_2$ esclusi per decisione del relatore; readout asimmetrico
opzionale, non prioritario; reti neurali (NQS) fuori scope. In più, per
il trimero: la topologia \textbf{catena aperta} resta un'estensione
possibile ma non prioritaria (decisione motivata e confermata dal
relatore, 6 settembre 2026) --- il risultato più originale della catena
(alternanza dei legami che cancella l'errore chirale di Trotter) resta
sul tavolo, non perso.

\section{Il filo conduttore verificato in ogni stadio}

Stessa proprietà usata nella pipeline del dimero, da riverificare qui
(non assumere dall'analogia): il readout simmetrico è un fattore
\textbf{moltiplicativo} sull'osservabile misurato, mai additivo --- da
cui, se confermato, discenderebbe che anche qui $N^*$ non dipende da
$p_\text{readout}$. Sul dimero questo era un risultato analitico
generale (non specifico a $N=2$); l'aspettativa è che si trasferisca
inalterato, ma resta da verificare esplicitamente allo stadio 5, non da
dare per scontato solo perché la dimostrazione originale non usava
dettagli specifici del numero di spin.

\end{document}
